\section*{第 \ref{ch:b2:metric} 章 --- 度量空间的拓扑}

\begin{solution}{exo:b2:metric:1}
$\delta$：对称性与分离性显然成立；三角不等式：
$\min(1, u + v) \leq \min(1,u) + \min(1,v)$，其中 $u, v \geq 0$（若
两个最小值之一为 $1$，则右边 $\geq 1$；否则为 $u + v$）。
$d$：它是 $\abs{\cdot}$ 经单射 $\arctan$ 的拉回：
三条公理随之成立。

收敛性：对 $\delta$，$\delta(x_n, x) \to 0 \iff \abs{x_n - x}
\to 0$（取值较小时两距离一致）：与通常的
收敛序列相同。对 $d$：$d(x_n, x) \to 0 \iff \arctan x_n \to
\arctan x \iff x_n \to x$（用到连续性与严格单调性，
$\arctan$ 及其反函数在相关取值范围上的性质）：仍
为通常的收敛。

$(\R, d)$ \emph{不}完备：$x_n = n$ 满足 $d(x_p, x_q)
= \abs{\arctan p - \arctan q} \to 0$（二者均趋于 $\frac\pi2$）：
是柯西列；但 $(x_n)$ 关于 $d$ 不收敛（若收敛，则其 $d$ 极限
须为普通极限，而该数列并无普通极限）。完备性是
\emph{距离}的性质，而不仅是收敛序列的性质。
\end{solution}

\begin{solution}{exo:b2:metric:2}
收敛 $\Rightarrow$ 柯西：$d(x_p, x_q) \leq d(x_p, \ell) +
d(\ell, x_q)$。有界性：超过 $N$ 后，$d(x_n, \ell) \leq 1$；其
前有限多项也都落在某个半径之内。

柯西 $+$ 收敛子列 $x_{\varphi(n)} \to \ell$：给定
$\varepsilon$，当 $n$ 充分大时，$d(x_n, \ell) \leq d(x_n,
x_{\varphi(n)}) + d(x_{\varphi(n)}, \ell) \leq 2\varepsilon$（其中
第一项由柯西性质控制，因 $\varphi(n) \geq n$）。

紧 $\Rightarrow$ 完备：柯西列必有
收敛子列（由紧性），从而收敛。
\end{solution}

\begin{solution}{exo:b2:metric:3}
$d_\infty(f, g) = \sup_{\intcc{0}{1}} \abs{x - x^2} = \frac14$
（$x - x^2$ 在 $x = \frac12$ 处取最大值）。

$\overline B(0, 1) = \{f : \sup\abs f \leq 1\}$：所有连续
函数，其取值在 $\intcc{-1}{1}$ 中。

$\{f : f(0) = 0\}$ 是 $\{0\}$ 在
\emph{求值}映射 $f \mapsto f(0)$ 下的原像，该映射是 $1$-利普希茨的
（$\abs{f(0) - g(0)} \leq d_\infty(f,g)$），从而连续：该集合
为闭集（由\cref{thm:b2:metric:globalcontinuity}可知）。同理 $\{f :
f(0) > 0\}$ 是开区间 $\intoo{0}{+\infty}$ 的原像，故为开集。
\end{solution}

\begin{solution}{exo:b2:metric:4}
\emph{完备：}柯西列在 $\varepsilon = \frac12$ 下
最终为常数列，从而收敛。

\emph{紧当且仅当有限：}若 $X$ 为有限集，则任一数列
必有某值出现无穷多次（即有常数子列）。若 $X$ 为无限集，
取两两互异的点列，其两两距离皆为 $1$：
则无柯西子列，更无收敛子列。

\emph{连通子集：}单点集（以及 $\emptyset$）。任一 $A$
若含两点 $x \neq y$，则可分解为 $\{x\} \cup (A
\setminus\{x\})$，二者在 $A$ 中皆开（离散空间的任一子集皆开 ---
半径为 $\frac12$ 的球皆为单点集）。
\end{solution}

\begin{solution}{exo:b2:metric:5}
函数 $g(x) = d(x, f(x))$ 在紧空间 $X$ 上
连续（由两次三角不等式可得 $\abs{g(x) - g(y)} \leq 2d(x,y)$），因此
它在某点 $a$ 处取得最小值
（由\cref{thm:b2:metric:compactprops}可知）。若 $f(a) \neq a$：
\[
g\bigl(f(a)\bigr) = d\bigl(f(a), f(f(a))\bigr) < d(a, f(a)) = g(a),
\]
与最小性矛盾。于是 $f(a) = a$；唯一性同样可证（若
有两个不动点 $a \neq b$，则 $d(a,b) = d(f(a), f(b)) < d(a,b)$）。

非紧的例子：$f(x) = x + \frac1x$，定义在 $\intco{1}{+\infty}$ 上：
$\abs{f(x) - f(y)} = \abs{x - y}\,\abs{1 - \frac{1}{xy}} < \abs{x -
y}$ 对 $x \neq y$ 成立（因 $xy > 1$），而 $f(x) > x$ 处处成立。
\end{solution}

\begin{solution}{exo:b2:metric:6}
任取 $x_n \in K_n$。自第 $n$ 项起，各项都属于 $K_n$；
特别地，整个数列落在紧集 $K_0$ 中，故可抽取
收敛子列 $x_{\varphi(k)} \to \ell$。对每个固定的 $n$，
$x_{\varphi(k)}$（其中 $\varphi(k) \geq n$）都落在\emph{闭}
集 $K_n$ 中，故极限 $\ell \in K_n$。因此 $\ell \in \bigcap K_n$。

闭集反例：$F_n = \intco{n}{+\infty}$（位于 $\R$ 中）：
它们嵌套、非空且为闭集，交集为空。
\end{solution}

\begin{solution}{exo:b2:metric:7}
$f^{-1}$ 连续，意味着：$f(F)$ 为闭集，其中 $F
\subseteq K$ 为闭集（关于 $f^{-1}$ 的原像即关于
$f$ 的像）。闭集 $F$ 位于紧集 $K$ 中，
故为紧集（由\cref{thm:b2:metric:compactprops}~(1)可知）；其连续像
$f(F)$ 是紧集，从而为闭集。于是 $f^{-1}$ 连续，即 $f$ 是
同胚。

反例：$\gamma(t) = (\cos t, \sin t)$，定义于
$\intco{0}{2\pi}$（非紧）上，它是映到单位圆上的
连续双射，但 $\gamma^{-1}$ 在 $(1,0)$ 处不连续：
单位圆上恰在坐标轴下方的点，其参数接近
$2\pi$，而非接近 $0$。
\end{solution}

\begin{solution}{exo:b2:metric:8}
$C = \bigcap_n C_n$，其中每个 $C_n$（由 $2^n$ 个
长度为 $3^{-n}$ 的闭区间之并）为闭集：$C$ 是闭且有界的，
在 $\R$ 中，因而为紧集（由\cref{thm:b2:metric:compactprops}~(2)可知）。

内部为空：$C$ 不含长度 $> 3^{-n}$ 的区间（它
包含于 $C_n$，而其各组成区间长度正是此值）
对每个 $n$ 成立。

基数：$C$ 的点恰为实数 $\sum_{n\geq1}
a_n 3^{-n}$，其中数字 $a_n \in \{0, 2\}$（每一步
删去的中间三分之一去掉数字 $1$）；映射 $(a_n) \mapsto
\sum a_n 3^{-n}$ 是从 $\{0,2\}^{\N^*}$ 到 $C$ 上的双射
（单射性同\cref{exo:b2:structures:3}）。故 $C$ 与
$\{0,1\}^{\N}$ 等势：不可数，但“长度为零”。
\end{solution}

\begin{solution}{exo:b2:metric:9}
假设 $a \notin f(X)$。像集 $f(X)$ 是紧的（连续
映射的像），因而为闭集；于是
\[
\varepsilon = d\bigl(a, f(X)\bigr)
= \inf_{y \in f(X)} d(a, y) > 0
\]
（紧集上连续函数的下确界必可取到；
若其为 $0$，则 $a$ 是闭集 $f(X)$ 的附着点，
从而属于它）。

考虑轨道 $x_n = f^n(a)$（$x_0 = a$）。对 $p < q$：
\[
d(x_p, x_q) = d\bigl(f^p(a), f^p(f^{q-p}(a))\bigr)
= d\bigl(a, f^{q-p}(a)\bigr) \geq \varepsilon,
\]
（等距映射迭代 $p$ 次；又 $f^{q-p}(a) \in f(X)$，因为 $q - p
\geq 1$）。两两距离均 $\geq \varepsilon$ 的点列没有
收敛子列 --- 与紧性矛盾。于是 $f(X) =
X$。
\end{solution}

\begin{solution}{exo:b2:metric:10}
设 $A, B \in GL_n(\C)$，并令 $p(z) = \det\bigl((1-z)A + zB\bigr)$：
关于 $z$ 的多项式（每个元素关于 $z$ 为仿射的；行列式是
关于各元素的多项式）。$p(0) = \det A \neq 0$：$p$ 不
恒为零，故它有有限多个根 $z_1, \dots, z_m$
（它们都不等于 $0$ 或 $1$：$p(1) = \det B \neq 0$）。
平面 $\C$ 去掉有限多个点后仍道路连通：从 $0$ 到 $1$ 必存在一条
避开各 $z_i$ 的道路（可取经过某点的折线，该点
远离所有根，或取圆弧；障碍只有有限多个）。
沿这样的道路 $\gamma$，$t \mapsto (1 - \gamma(t))A +
\gamma(t)B$ 是一条连续道路，且\emph{完全位于} $GL_n(\C)$ 之中，从 $A$
到 $B$（行列式在其上处处不为零）。因此 $GL_n(\C)$
道路连通 --- 不同于其实数版本
（见\cref{ex:b2:metric:glnr}）：复平面有足够空间
绕开障碍。
\end{solution}

\begin{solution}{exo:b2:metric:11}
\emph{利普希茨：}对 $a \in A$，$d(x, a) \leq d(x, y) + d(y,
a)$。对 $a$ 取下确界得 $d(x, A) \leq d(x, y) + d(y,
A)$，交换 $x, y$，即得 $\abs{d(x,A) - d(y,A)} \leq d(x,y)$。

\emph{零值情形：}$d(x, A) = 0$ 当且仅当存在 $a_n \in A$，使
$d(x, a_n) \to 0$，当且仅当 $x$ 为 $A$ 中点的极限，当且仅当 $x \in
\overline A$。

\emph{切换构造：}对不交的闭集 $A, B$，分母
$d(x,A) + d(x,B)$ 永不为零（否则会迫使 $x \in \overline
A \cap \overline B = A \cap B = \emptyset$），故 $\varphi$ 是良
定义的，且作为连续函数之商而连续
（分母处处不为零）。$\varphi(x) = 0$ 当且仅当 $d(x, A) = 0$
即 $x \in A$；$\varphi(x) = 1$ 当且仅当 $d(x, B) = 0$，即 $x \in B$；
且处处有 $0 \leq \varphi \leq 1$。
\end{solution}

\begin{solution}{exo:b2:metric:12}
设 $B(x_0, r_0)$ 为任一球，则在其中可找到 $\bigcap U_n$
的一点。由于 $U_1$ 稠密且开，$U_1 \cap B(x_0, r_0)$ 是
非空开集：它包含一个闭球 $\overline B(x_1,
r_1)$，满足 $0 < r_1 \leq \frac{r_0}2$（缩小半径）。
归纳地，$U_{n+1} \cap B(x_n, r_n)$ 非空且开：取
$\overline B(x_{n+1}, r_{n+1}) \subseteq U_{n+1} \cap B(x_n,
r_n)$，使 $r_{n+1} \leq \frac{r_n}2$。对 $p, q \geq n$，
$x_p, x_q$ 都位于 $B(x_n, r_n)$ 中，且 $r_n \leq 2^{-n}r_0$：
点列为柯西列，于是收敛到某个 $\ell$，
（由完备性）。对每个 $n$，点列的尾部位于
\emph{闭}球 $\overline B(x_{n+1}, r_{n+1}) \subseteq U_{n+1}
\cap B(x_0, r_0)$，故 $\ell \in U_{n+1}$（对每个 $n$），且 $\ell
\in \overline B(x_1, r_1) \subseteq B(x_0, r_0)$，因此
$\bigcap_n U_n$ 与每个球都相交，故稠密。

\emph{应用：}若 $\R = \bigcup_n F_n$，其中 $F_n$ 为闭集
且内部为空，则 $U_n = \R \setminus F_n$ 稠密
（$\overline{U_n} = \R$ 当且仅当 $F_n$ 的内部为空），且为开集，
再由贝尔纲定理，在 $\bigcap U_n = \R \setminus \bigcup
F_n$ 中存在一点：矛盾。特别地，$\R \neq \bigcup_{x \in D}
\{x\}$，其中 $D$ 为可数集（单点集为闭集，其内部为空
）。于是 $\R$ 不可数 ---
\cref{thm:b2:structures:cantor}的另一种证法。
\end{solution}

\begin{solution}{pb:b2:metric:1}
\textbf{1.} 设 $F$ 为完备度量空间 $X$ 中的闭集，$(x_n)
\subseteq F$ 为柯西列：它在 $X$ 中收敛到某个 $\ell$，
$\ell \in F$，因为 $F$ 是闭的（$F$ 中点列的极限
仍在 $\overline F = F$ 中）。于是 $F$ 完备。反之，设 $A
\subseteq X$ 完备，$x \in \overline A$。
$A$ 中存在点列收敛到 $x$；它为柯西列，故它\emph{在}
$A$ 中收敛，极限唯一，故 $x \in A$：$A$ 为闭集。由于
$\bigl(C(I), d_\infty\bigr)$ 完备（
\cref{thm:b2:metric:rncomplete}），其闭子集
都完备。

\textbf{2.} 对 $y, z \in C(I)$ 及 $t \in I$：
\[
\Bigl|\int_{t_0}^{t}y - \int_{t_0}^{t}z\Bigr| \leq
\abs{t - t_0}\,\sup_I\abs{y - z} \leq h\,d_\infty(y, z),
\]
再对 $t$ 取上确界。

\textbf{3.} 利用两个不动点方程与三角
不等式：
\[
d(\ell_g, \ell_{\widetilde g})
= d\bigl(g(\ell_g), \widetilde g(\ell_{\widetilde g})\bigr)
\leq d\bigl(g(\ell_g), g(\ell_{\widetilde g})\bigr) +
d\bigl(g(\ell_{\widetilde g}), \widetilde g(\ell_{\widetilde
g})\bigr)
\leq k\,d(\ell_g, \ell_{\widetilde g}) +
d\bigl(g(\ell_{\widetilde g}), \widetilde g(\ell_{\widetilde
g})\bigr),
\]
解出 $d(\ell_g, \ell_{\widetilde g})$（系数
$1 - k$ 为正）。第二个不等式将求值后的差
以一致差为界。

\textbf{4.} $g^m$ 是非空完备空间上的压缩映射：
它有唯一不动点 $\ell$
（由\cref{thm:b2:metric:banach}）。于是 $g^m(g(\ell)) =
g(g^m(\ell)) = g(\ell)$。$g(\ell)$ 是 $g^m$ 的不动点，
故 $g(\ell) = \ell$（由唯一性）。$g$ 的任一不动点都是
$g^m$ 的不动点，这即 $g$ 的唯一性。轨道：取定 $r \in \{0, \dots,
m-1\}$；子列 $(x_{qm + r})_q$ 是 $g^m$ 的轨道，
起点为 $x_r$，故收敛于 $\ell$（$q \to \infty$ 时，
再用巴拿赫定理）。$m$ 个子列都收敛于同一个
$\ell$，从而 $x_n \to \ell$：给定 $\varepsilon$，每个剩余类
从某项起都落在 $\varepsilon$ 之内，
而剩余类只有有限多个。

\textbf{5.} 若 $y$ 连续且取值于 $\intcc{y_0 -
b}{y_0 + b}$ 中，则被积函数 $s \mapsto f(s, y(s))$
在 $I$ 上连续（复合而得），因此右端项为 $C^1$，
导数为 $f(t, y(t))$（微积分基本定理，
参见本科一年级分册）。若 $y$ 满足积分方程，则它就是该
$C^1$ 函数，$y(t_0) = y_0$，且 $y' = f(t, y)$。反之，
将 $y' = f(s, y(s))$ 从 $t_0$ 到 $t$ 积分，
即得积分方程。

\textbf{6.} $X_h$ 含常值函数 $y_0$；它是闭集，
因为它是 $\intcc{0}{b}$ 在连续映射 $y
\mapsto d_\infty(y, y_0)$ 下的原像（距离函数为 $1$-利普希茨），因而
由第 1 题它完备。稳定性：对 $y \in X_h$ 与 $t \in
I$，
\[
\abs{T(y)(t) - y_0} = \Bigl|\int_{t_0}^{t}f(s, y(s))\dd s\Bigr|
\leq M\abs{t - t_0} \leq Mh \leq b ,
\]
末一步由 $h \leq b/M$ 得出（若 $M = 0$，则平凡）。又 $T(y)$
连续（甚至为 $C^1$，由第 5 题），故 $T(y) \in X_h$。

\textbf{7.} 对 $t \in I$：
\[
\abs{T(y)(t) - T(z)(t)} \leq \int_{t_0}^{t}\abs{f(s, y(s)) -
f(s, z(s))}\,\abs{\dd s} \leq L\abs{t - t_0}\,d_\infty(y,z)
\leq Lh\,d_\infty(y,z).
\]
若 $Lh < 1$，则 $T$ 是非空完备
空间 $X_h$ 上的压缩映射，由巴拿赫定理有唯一不动点 ---
据第 5 题，即唯一解。

\textbf{8.} 用归纳法；$n = 1$ 的情形即
第 7 题的中间不等式。设对 $n$ 的估计成立，取 $t \geq t_0$（
$t \leq t_0$ 的情形对称）：
\[
\abs{T^{n+1}(y)(t) - T^{n+1}(z)(t)}
\leq L\int_{t_0}^{t}\abs{T^n(y)(s) - T^n(z)(s)}\dd s
\leq L\int_{t_0}^{t}\frac{L^n(s - t_0)^n}{n!}\dd s\;
d_\infty(y,z),
\]
积分值为 $\frac{L^n(t -
t_0)^{n+1}}{(n+1)!}$，即对 $n + 1$ 成立的估计。于是
$d_\infty(T^n y, T^n z) \leq \frac{(Lh)^n}{n!}d_\infty(y, z)$，
而 $\frac{(Lh)^n}{n!} \to 0$（指数级数
收敛）。故某个 $T^m$ 为压缩映射。对完备的 $X_h$ 用第 4 题，
即 $T$ 有唯一不动点，亦即柯西问题在 $I$ 上
恰有一个取值于
$\intcc{y_0 - b}{y_0 + b}$ 中的解。

\textbf{9.} 设 $y$ 是 $I$ 上问题的解，并设集合
$E = \{t \in I, t > t_0 : \abs{y(t) - y_0} > b\}$ 非空
（$t < t_0$ 一侧对称）；令 $\tau = \inf E$。由
连续性，$\abs{y(s) - y_0} \leq b$ 对 $s \in
\intcc{t_0}{\tau}$ 成立，故图像在该段位于 $R$ 中，积分
方程在 $\tau$ 之前成立，且
\[
\abs{y(\tau) - y_0} = \Bigl|\int_{t_0}^{\tau}f\bigl(s,
y(s)\bigr)\dd s\Bigr| \leq M(\tau - t_0) \leq Mh \leq b .
\]
若 $\tau < t_0 + h$，则 $E$ 中从右侧任意接近 $\tau$
的点由连续性给出 $\abs{y(\tau) - y_0} \geq
b$，从而 $= b$；但此时上式迫使 $M(\tau - t_0) =
Mh$，即 $\tau = t_0 + h$，矛盾。故 $\tau = t_0 + h$，
$E \subseteq \{t_0 + h\}$，且上式（取 $\tau = t_0 + h$）
给出 $\abs{y(t_0 + h) - y_0} \leq b$，与
$E$ 的定义相矛盾。故 $E = \emptyset$：$I$ 上每个解都留在带中，
是 $T$ 在 $X_h$ 中的不动点，唯一性
无条件成立。

\textbf{10.} $T(y)(t) = 1 + \int_0^t y$。从 $y^{(0)} \equiv
1$：
\[
y^{(1)}(t) = 1 + t,\quad
y^{(2)}(t) = 1 + t + \frac{t^2}2,\quad\dots\quad
y^{(n)}(t) = \sum_{k=0}^{n}\frac{t^k}{k!}
\]
（归纳法：对部分和积分即加上下一项）。
它们是 $\eu^{\,t}$ 的泰勒部分和；在任意有界的
区间 $I$ 上它们一致收敛于 $\eu^{\,t}$（尾部
被收敛的数项级数 $\sum h^k/k!$ 所控制），
而它正是唯一解。

\textbf{11.} $y \equiv 0$ 是一个解。对 $y_c$：它是 $C^1$
（两段都是，且在 $t = c$ 处导数匹配：$0$ 与
$2(t - c) \to 0$），当 $t > c$ 时：$y_c' = 2(t - c) =
2\sqrt{(t-c)^2} = 2\sqrt{\abs{y_c}}$；当 $t \leq c$ 时，
两边都为零。于是初值问题 $y(0) = 0$ 有无穷多个
解（$c \geq 0$ 任意，以及 $y \equiv 0$）。右端函数
$\varphi(y) = 2\sqrt{\abs y}$ 在 $0$ 附近不满足利普希茨条件：
$\frac{\varphi(y) - \varphi(0)}{y - 0} = \frac{2}{\sqrt y} \to
+\infty$，当 $y \to 0^+$ 时：不存在常数 $L$，使其在
$0$ 的任何邻域上成立 --- 这正是所有解
分岔之处。

\textbf{12.} 分离变量（或直接验证）可得唯一的局部解
为 $y(t) = \frac1{1 - t}$，其定义于
$\intoo{-\infty}{1}$，并在 $t = 1$ 处爆破。对于
矩形 $\intcc{-a}{a} \times \intcc{1 - b}{1 + b}$：$M =
\sup y^2 = (1 + b)^2$，于是可保证的半宽为 $h =
\min\bigl(a, \frac{b}{(1+b)^2}\bigr)$。求
$\frac{b}{(1+b)^2}$ 的最大值：导数在 $b = 1$ 处为零，此时值为
$\frac14$。因此该定理只保证解在
$\intcc{-\frac14}{\frac14}$ 上存在 --- 这正确地短于真实
存在区间：向前为 $1$，向后则差无穷多。该定理
本质上是局部的，爆破表明它不可能是其他情形。

\textbf{13.} $g$ 把 $\intcc12$ 映入自身：$g$ 在
$\intcc{1}{\sqrt2}$ 上递减，此后递增（考察 $g'(x) = \frac12
- \frac1{x^2}$），且 $g(1) = g(2) = \frac32$，最小值
为 $g(\sqrt2) = \sqrt2 > 1$：$g(\intcc12) \subseteq
\intcc{\sqrt2}{\frac32} \subseteq \intcc12$；又 $g$ 把
有理数映为有理数。压缩性：$\abs{g'(x)} = \abs{\frac12
- \frac1{x^2}} \leq \frac12$ 在 $\intcc12$ 上成立（$\frac1{x^2} \in
\intcc{\frac14}{1}$），由中值不等式得
$\abs{g(x) - g(y)} \leq \frac12\abs{x - y}$。不动点
满足 $\frac x2 = \frac1x$，即 $x^2 = 2$，这在
$\Q$ 中不可能。失效的假设是 $X$ 的完备性（$\Q \cap
\intcc12$ 不完备）；在其完备化 $\intcc12$ 中，
不动点为 $\sqrt2$ --- 巴拿赫不动点定理在有理数上运行
时\emph{创造出}了无理数。

\textbf{14.} 后验估计：$d(x_n, \ell) \leq d(x_n, x_{n+1}) +
d(x_{n+1}, \ell) \leq k\,d(x_{n-1}, x_n) + k\,d(x_n, \ell)$,
由此 $d(x_n, \ell) \leq \frac{k}{1-k}d(x_n, x_{n-1})$。海伦迭代
自 $x_0 = \frac32$ 出发：$x_1 = \frac{17}{12}$，$d(x_1, x_0) =
\frac1{12}$，$k = \frac12$：先验界
$\frac{k^n}{1-k}d(x_1,x_0) = \frac{2^{-n+1}}{12}$ 首次降至
$10^{-6}$ 以下时 $n = 18$。实际上 $x_1 = \frac{17}{12}
\approx 1.41667$（误差 $2.5\cdot10^{-3}$），$x_2 =
\frac{577}{408} \approx 1.4142157$（误差 $2.1\cdot10^{-6}$），
$x_3 \approx 1.41421356237469$（误差 $1.6\cdot10^{-12}$）：
三步即足够。每一步海伦迭代大致把
误差\emph{平方}（二次收敛，一种牛顿现象：
见\cref{ch:b2:realfun}）；压缩估计仅
把误差减半，对最坏情形是诚实的，但在此处偏于保守。

\textbf{15.} 取第 3 题中的 $g = T_{y_0}$（一个
$Lh$ 压缩，$Lh < 1$），$\widetilde g = T_{z_0}$，其
不动点为 $y[z_0]$。对任意 $y$，
\[
\abs{T_{y_0}(y)(t) - T_{z_0}(y)(t)} = \abs{y_0 - z_0},
\]
（两个积分完全相同），故 $\sup_y
d_\infty(T_{y_0}(y), T_{z_0}(y)) = \abs{y_0 - z_0}$，因此
由第 3 题得
\[
d_\infty(y[y_0], y[z_0]) \leq \frac{\abs{y_0 - z_0}}{1 - Lh} .
\]

\textbf{16.} 在每一段上应用第 15 题，
取左端点的值作初值，界定
右端点处的偏差：
\[
d_\infty \leq \frac{1}{1 - 1/2}\,(\text{左端偏差})
= 2\,(\text{左端偏差}) .
\] 对
$m$ 段作归纳，末端偏差至多为
$2^m\abs{y_0 - z_0}$，而整个
区间上的一致偏差满足同一界（每段的上确界在其阶段
受控）。取段长 $h \asymp \frac1{2L}$，则
因子为 $2^m = 2^{\,\text{长度}\cdot 2L}$：关于
区间的长度呈指数增长，恰如格朗沃尔界
$\eu^{L\abs{t-t_0}}$ 所预示的那样，但常数更好。

\textbf{17.} 方案相同：对 $y \in X_h$，
\[
\abs{T_f(y)(t) - T_g(y)(t)} \leq \int_{t_0}^t \abs{f(s,y(s)) -
g(s,y(s))}\,\abs{\dd s} \leq \varepsilon h ,
\]
于是由第 3 题（$T_f$ 为压缩，$T_g$ 为
扰动映射）得 $d_\infty \leq \frac{\varepsilon h}{1 - Lh}$。

\textbf{18.} 由第 17 题，取 $f = f_\lambda$，$g =
f_\mu$：$d_\infty(y_\lambda, y_\mu) \leq \frac{Ch}{1 -
Lh}\abs{\lambda - \mu}$：解映射是利普希茨的，
其常数为 $\frac{Ch}{1-Lh}$。

\textbf{19.} 所用论证只涉及：度量公理、
该阶段的完备性、界 $\abs{\int} \leq
\int\abs{\cdot}$，以及 $f$ 的利普希茨性质 --- 这些
对 $\R^n$ 值函数均成立，其中 $d_\infty$ 建立在
\cref{ex:b2:metric:examples}中的任一等价距离之上
（由\cref{thm:b2:metric:rncomplete}可知其完备）。对 $y' = Ay$，
其中 $A$ 为幂零矩阵：$y^{(0)} \equiv (c_1, c_2)$，
\[
y^{(1)}(t) = (c_1, c_2) + \int_0^t (c_2, 0)\,\dd s
= (c_1 + tc_2,\; c_2),
\]
又 $Ay^{(1)}(s) = (c_2, 0)$ 再次成立，故 $y^{(2)} = y^{(1)}$。
迭代从 $n = 1$ 起保持不变，停在精确解
$y(t) = (c_1 + tc_2, c_2) = \eu^{tA}y(0)$ --- 幂零性
截断了指数级数，而皮卡迭代察觉到了这一点。

\textbf{20.} 固定 $y \in X$，令 $g_y(x) = y - \eta(x)$，它是
$k$-压缩映射，定义于完备的 $X$ 上，故恰有一个 $x$
使 $x + \eta(x) = y$，即映射 $\Phi = \mathrm{id} + \eta$ 为
双射。逆映射的利普希茨性：若 $\Phi(x) = y$ 且 $\Phi(x') =
y'$,
\[
d(x, x') \leq d(y, y') + d\bigl(\eta(x), \eta(x')\bigr)
\leq d(y, y') + k\,d(x, x'),
\]
于是 $d(x, x') \leq \frac{1}{1-k}d(y, y')$。（此处的距离
来自范数结构，故 $d(a - c, b - c) = d(a, b)$，
这正是第一个不等式所用到的。）

\textbf{21.} $g(x) = m + e\sin x$ 是 $e$-利普希茨的，
在完备的 $\R$ 上（由中值不等式，$\abs{g'} = \abs{e\cos x}
\leq e < 1$）：由巴拿赫不动点定理，唯一解存在，
且迭代全局收敛。取 $e = \frac12$，$m = 1$，$x_0
= 1$：$x_1 = 1 + \frac12\sin 1 \approx 1.42074$，$d(x_1, x_0)
\approx 0.4207$，而先验界
$\frac{(1/2)^n}{1/2}\cdot 0.4207 \leq 10^{-3}$ 首次成立于
$n = 10$：十次迭代即得到确认（真值 $x \approx
1.4987$ 实际上在 $10^{-3}$ 的精度内即已达到，约在 $n =
5$）。

\textbf{22.} 利用数字描述
（见\cref{exo:b2:metric:8}）：$C$ 是所有这类和之集：
$\sum_{n\geq1}a_n3^{-n}$，$a_n \in \{0, 2\}$。于是 $S_1(C) =
\{x/3 : x \in C\}$ 为 $a_1 = 0$ 的子集，而 $S_2(C) =
\{x/3 + 2/3\}$ 为 $a_1 = 2$ 的子集：二者之并，
在 $a_1$ 自由选取时恰为 $C$。唯一性一句话即可说明：
在非空紧子集全体所成的空间上
（$\intcc01$ 的子集，以豪斯多夫距离为度量），$A \mapsto
S_1(A) \cup S_2(A)$ 是一个 $\frac13$-压缩映射，作用于完备的
空间上，故由巴拿赫不动点定理只允许一个不动集 --- 本科
三年级分册使豪斯多夫度量与此论证严格化。

\textbf{23.} 令 $Z = \{t \in J : y(t) = z(t)\}$：非空
（$t_0 \in Z$），在 $J$ 中是闭集（两连续映射的相等点集，即
$\{0\}$ 在 $y - z$ 下的原像）。开性：若 $t_1 \in Z$，应用
局部定理（问题 8）于点 $(t_1, y(t_1))$，
在使 $f$ 为利普希茨的矩形内：于一小区间上
$t_1$ 附近，$y$ 与 $z$ 都解同一柯西问题，故
二者在该区间上重合（即问题 9 的无条件唯一性）：
$t_1$ 的一个邻域含于 $Z$。
区间 $J$ 的非空子集，若在 $J$ 中既开又闭，则必为整个 $J$
（区间连通，见
\cref{thm:b2:metric:connectedness}）：$y = z$ 于 $J$ 上。

\textbf{24.} 这里 $f(t, y) = \alpha(t)y + \beta(t)$ 是
$L$-利普希茨的（关于 $y$，在整个 $\intcc AB \times \R$ 上），其中 $L =
\sup\abs\alpha$（有限：$\alpha$ 在线段上连续），且
无需任何带形区域 $\intcc{y_0 - b}{y_0 + b}$：取 $X = C(\intcc
AB)$ 整体，$T$ 在其上良定义。问题 8 中的归纳法
原样成立，给出 $d_\infty(T^ny, T^nz) \leq
\frac{(L(B - A))^n}{n!}d_\infty(y, z)$：某次迭代是
压缩映射（无论长度 $B - A$ 为何），由问题 4
可得整个线段上恰有一个解。
线性性恰好起作用一次：它使利普希茨界
对 $y$ 全局成立，去掉了矩形及其 $h \leq b/M$。

\textbf{25.} \emph{完备性}使迭代所得的柯西列
变为真正的解（问题 1、6、8），而其缺失
则使 $\sqrt2$ 能逃出 $\Q$（问题 13）。
\emph{压缩}给出唯一性、算法及
误差界（问题 7、14）。\emph{迭代技巧}去除了
小量条件 $Lh < 1$，使所确认的区间
仅依赖于 $M$，而不依赖于 $L$ --- 并使线性方程
成为全局的（问题 8、24）。\emph{连通性}将局部
唯一性提升为全局唯一性（问题 23）。
反例中，$2\sqrt{\abs y}$ 说明利普希茨条件不可缺（问题
11），$y^2$ 说明局部性不可缺（问题 12），$\Q$ 说明
完备性不可缺（问题 13）。最高峰为
皮卡--林德勒夫定理（问题 8）；当 $f$ 仅
为连续时，存在性依然成立，唯一性则失效，而
证明以函数集的紧性取代压缩性，
--- 皮亚诺定理借助阿尔泽拉--阿斯科利定理，在
本科三年级分册中给出。
\end{solution}
